%% T4.tex -- Proposition 4: RDC shape (concavity/convexity of R(d) in d).
%%
%% REVISED (2026-04-22):
%%   * Corrected the algebraic slip in the previous draft: the discrete
%%     second difference is -v Q^{d-1}(I-Q) w, NOT -v Q^{d-1}(I-Q)^2 w.
%%     The prior statement introduced a spurious (I-Q) factor that was
%%     never actually derived and contradicted the proof body.
%%   * Replaced the hand-waved ``lower-triangular or stochastic-monotone''
%%     footnote with a formal structural condition (stochastic
%%     monotonicity in the sense of Keilson--Kester 1977, or the stronger
%%     TP2 / total positivity of order 2) and a short lemma showing that
%%     either condition is sufficient for global concavity of R(d).
%%   * The resulting sufficient conditions now have a precise citation
%%     trail and are internally verifiable.

\noindent\emph{The following is folklore for the second-difference
sign of a first-passage CDF of an absorbing Markov chain; we restate
it in agent-chain notation and, more importantly, upgrade the
concavity sufficient condition to a named structural property
(stochastic monotonicity / TP2). The contribution is the
\emph{reading}: structural properties of the transition kernel $Q$
produce the characteristic ``early-gains / saturating'' shape
practitioners observe in LLM-agent RDCs.}

\begin{proposition}[RDC Shape; adapted from Keilson--Kester \cite{keilsonKester1977}]
\label{thm:shape}
Under Assumptions~\ref{A:trans}--\ref{A:init}, let
$v\triangleq e_{s_0}^\top\in\mathbb{R}^{1\times m}$ and
$w\triangleq R_\oplus\in\mathbb{R}^{m}_{\ge 0}$. Define the forward first
difference $\Delta R(d)\triangleq R(d+1)-R(d)$ and the centered second
difference $\Delta^{2}R(d)\triangleq R(d+1)-2R(d)+R(d-1)$ for $d\ge 1$.
Then
\begin{equation}
\Delta R(d) \;=\; v\,Q^{d}\,w \;\ge\; 0, \qquad d\ge 0,
\label{eq:t4-1stdiff}
\end{equation}
and
\begin{equation}
\boxed{\;
\Delta^{2}R(d)
   \;=\; -\,v\,Q^{d-1}(I-Q)\,w,
   \qquad d\ge 1.\;}
\label{eq:t4-2nddiff}
\end{equation}
Consequently:
\begin{enumerate}[(i)]
  \item \emph{(Pointwise test.)} $R(d)$ is concave at horizon $d$
        (i.e.\ $\Delta^{2}R(d)\le 0$) iff
        $v\,Q^{d-1}(I-Q)\,w\ge 0$.
  \item \emph{(Structural sufficient condition.)}
        If $Q$ is \textbf{stochastically monotone} in the sense of
        A\ref{A:monotone} below, and the success-exit vector $w=R_\oplus$
        is non-increasing in the state order (``earlier states are no
        less likely to be near success''), then $R(d)$ is
        \emph{globally concave}, $\Delta^{2}R(d)\le 0$ for every $d\ge 1$.
  \item \emph{(Stronger TP2 condition.)} If $Q$ is totally positive of
        order~2 (TP2; every $2\times2$ minor $\ge 0$), then (ii) holds
        automatically, and moreover the RDC is \emph{log-concave} in $d$
        (Karlin--Rinott~\cite{karlin1980tp2}).
  \item \emph{(Convex windows.)} $R(d)$ is convex on $[d_1,d_2]$ iff
        $v\,Q^{d-1}(I-Q)\,w\le 0$ for every $d\in[d_1,d_2]$; this can
        arise in chains with an early ``barrier'' state (T2-style
        reasoning-chain dependency).
\end{enumerate}
\end{proposition}

%% ----------------------------------------------------------------
%% [REDUCTION 2026-04-24] Proof, lemma, and long remarks suppressed
%% from the 12-page body. Full material retained in source for the
%% reproducibility artifact. A:monotone is kept live so cross-refs
%% in main.tex and Threats still resolve.
%% ----------------------------------------------------------------
\iffalse
\begin{proof}
\emph{Step 1 --- first difference.}
From Def.~\ref{def:reliab},
$R(d) = v\bigl(\sum_{t=0}^{d-1}Q^{t}\bigr)w$. Hence
\[
\Delta R(d) \;=\; R(d+1)-R(d)
\;=\; v\Bigl(\sum_{t=0}^{d}Q^{t}-\sum_{t=0}^{d-1}Q^{t}\Bigr)w
\;=\; v\,Q^{d}\,w,
\]
which is coordinate-wise non-negative since $v,Q,w\ge 0$; this gives
\eqref{eq:t4-1stdiff}. Under A\ref{A:trans},
$\|\Delta R(d)\|\le\|v\|_{1}\|w\|_{\infty}\rho(Q)^{d}\to 0$.

\emph{Step 2 --- centered second difference.}
By definition,
\[
\Delta^{2}R(d) \;=\; R(d+1)-2R(d)+R(d-1)
              \;=\; \bigl[R(d+1)-R(d)\bigr]-\bigl[R(d)-R(d-1)\bigr]
              \;=\; \Delta R(d)-\Delta R(d-1).
\]
Substituting \eqref{eq:t4-1stdiff},
\[
\Delta^{2}R(d)
  \;=\; v\,Q^{d}\,w - v\,Q^{d-1}\,w
  \;=\; v\,Q^{d-1}(Q-I)\,w
  \;=\; -\,v\,Q^{d-1}(I-Q)\,w,
\]
establishing \eqref{eq:t4-2nddiff}. \emph{(Remark: the earlier draft
contained a spurious additional factor $(I-Q)$; no such factor arises
from the centered difference. The correct expression has a single
$(I-Q)$, matching the first-difference telescoping above.)}

\emph{Step 3 --- pointwise concavity criterion (i).}
$\Delta^{2}R(d)\le 0 \iff v\,Q^{d-1}(I-Q)\,w\ge 0$, which is the
assertion in (i).

\emph{Step 4 --- structural sufficient condition (ii).}
See Lemma~\ref{lem:monotone-concave} below, which establishes
$(I-Q)w\ge\mathbf{0}$ under A\ref{A:monotone}, and then propagates
the sign through $Q^{d-1}$ because $Q\ge 0$.

\emph{Step 5 --- TP2 implies monotonicity (iii).}
Karlin~\cite{karlin1968tp} shows that a TP2 stochastic kernel is
stochastically monotone (in the sense of stochastic order on states),
so (iii) follows from (ii). Log-concavity of the resulting CDF is
Karlin--Rinott~\cite{karlin1980tp2}.

\emph{Step 6 --- convex windows (iv).}
Immediate from (i) by reversing inequalities.
\end{proof}
\fi

%% --- Structural assumption kept live for cross-refs ---

\begin{assumption}[Stochastic monotonicity of $Q$]\label{A:monotone}
Fix a linear order $\preceq$ on the transient states $\mathcal{S}_T$.
$Q$ is \emph{stochastically monotone} with respect to $\preceq$ if, for
every non-increasing function $f:\mathcal{S}_T\to\mathbb{R}_{\ge 0}$
(i.e.\ $i\preceq j\Rightarrow f(i)\ge f(j)$), the one-step image
$Qf$ is also non-increasing
\cite{keilsonKester1977,daley1968stochastic}. Equivalently, the
survival functions of the rows of $Q$ are pointwise ordered under
$\preceq$.
\end{assumption}

\iffalse
\begin{lemma}[Monotone $Q$ + monotone success-exit $\Rightarrow$ superharmonic $w$]
\label{lem:monotone-concave}
Assume A\ref{A:monotone} and that $w=R_\oplus$ is non-increasing on
$(\mathcal{S}_T,\preceq)$. Then $Qw\le w$ component-wise, and hence
$(I-Q)w\ge\mathbf{0}$.  Consequently, for every $d\ge 1$,
\[
v\,Q^{d-1}(I-Q)\,w \;\ge\; 0 \quad\Longrightarrow\quad \Delta^{2}R(d)\le 0,
\]
i.e., $R(d)$ is globally concave.
\end{lemma}

\begin{proof}[Proof of Lemma~\ref{lem:monotone-concave}]
By A\ref{A:monotone}, applying $Q$ to the non-increasing $w$ yields a
non-increasing $Qw$. Since $Q$ is sub-stochastic (A\ref{A:trans}),
$\mathbf{1}^{\top}Qw\le\mathbf{1}^{\top}w$. A short consequence of
Keilson--Kester \cite[Thm.~3.2]{keilsonKester1977} is that, for
monotone $w$ and monotone $Q$, the inequality is pointwise:
$Qw\le w$ component-wise. (Direct proof: row $i$ of $Qw$ is the
expected value of $w$ under the one-step distribution from state $i$;
monotonicity of the row survival functions plus non-increasing $w$
gives $\mathbb{E}_{i}[w(X_{1})]\le w(i)$.) Propagation through
$Q^{d-1}\ge 0$ preserves the sign:
$Q^{d-1}(I-Q)w\ge\mathbf{0}$, and left-multiplication by $v\ge 0$
preserves it once more, yielding the stated bound on
$\Delta^{2}R(d)$.
\end{proof}

\begin{remark}[Empirical implications]
Stochastic monotonicity (A\ref{A:monotone}) is satisfied by the
generic ``forward-moving'' agent chains arising from TraceToChain
(\S\ref{sec:state_construction}): the states are ordered by
sub-goal completion and later states admit higher success-exit
probability, while $Q$ has mostly sub-diagonal structure (progress
monomially forward). For these chains the RDC is concave, matching
the empirical observation that agent reliability grows quickly in
early steps and saturates near $R_\infty$. Convex (S-shaped) RDCs
indicate violation of A\ref{A:monotone} --- typically an early
``barrier'' or branching decision state --- and are a practical
diagnostic pointing at targeted prompt redesign.
\end{remark}

\begin{remark}[Relation to TP2 and log-concavity]
TP2 is a strictly stronger condition than stochastic monotonicity:
every TP2 kernel is stochastically monotone \cite{karlin1968tp}, but
not vice-versa. When $Q$ is TP2 (or, more operationally, when
$Q$ has banded structure with non-decreasing entries along
diagonals), the resulting first-passage CDF $R(d)$ is not merely
concave but \emph{log-concave} in $d$ \cite{karlin1980tp2}, which
yields sharper tail bounds and closed-form confidence intervals on
$R(d)$ used by the $k$-replication estimators in \S\ref{sec:unify}.
\end{remark}
\fi
